\documentclass{ximera}
\graphicspath{ {./images/} }

\title{MATH 1075 }

\author{}
\date{}


\begin{document}
Print Name: $\_\_\_\_$\\
OSU Name.\#: $\_\_\_\_$\\
Instructor: $\_\_\_\_$\\
Signature: $\_\_\_\_$

\section*{Midterm 1}
\section*{Autumn 2019}
\section*{Instructions:}
\begin{itemize}
  \item You have 55 minutes to complete this exam.
  \item Circle or box all final answers.
  \item A calculator may be used given the calculator policy outlined in the syllabus. Even if a calculator is used, you must show all work on each problem to receive full credit.
  \item Write clearly and legibly to receive credit.
  \item Do not round; give only exact answers (unless specified otherwise).
\end{itemize}

Problem 1. Solve the inequality. Show your work. Your final answer should be simplified and in interval notation. Graph the solution set to the inequality on the number line below.\\
(10 pts)

$$-7(X+5)+32 \leq 3(7-X)$$

\begin{center}
\includegraphics[alt={Number line ranging from negative 10 to 10},width=\textwidth]{number-line.jpg}
\end{center}

Problem 2. Solve the inequality. Show your work. Your final answer should be simplified and in interval notation.

$$-15<2 R-5 \leq-1$$

Problem 3. Solve for $W$. Show your work. If there is no solution, write "No solution." (6 pts)

$$|8 W+8|-5=11$$

Problem 4. Solve the inequality. Show your work. Your final answer should be simplified and in interval notation. If there is no solution, write "No solution." ( 8 pts)

$$3|Q+2|-5>-4$$

Problem 5. Factor the expression completely.

$$30 g^{5} y^{8}+5 g^{9} y^{7} q^{3}$$

Problem 6. Factor the expression completely.

$$3 A B-5 A+3 B C-5 C$$

Problem 7. Factor the expression completely.

$$m^{2}+6 m-16$$

Problem 8. Factor the expression completely.

$$48 w^{2}-152 w+80$$

Problem 9. Factor the expression completely.

$$-18 V^{2}-60 V-50$$

Problem 10. Factor the expression completely.\\
$9-25 h^{2}$

Problem 11. Factor the expression completely.

$$z^{4}-81$$

Problem 12. Factor the expression completely. You may find the following formula to be helpful: $a^{3}-b^{3}=(a-b)\left(a^{2}+a b+b^{2}\right)$.\\
(8 pts)\\
$64-27 w^{3}$

Problem 13. True or False?\\
(8 pts)\\
(a) The expression $(a+b)^{2}$ can always be rewritten as $a^{2}+b^{2}$.

True False\\
(b) The expression $(a-b)^{2}$ can always be rewritten as $a^{2}-2 a b+b^{2}$.

True False\\
(c) If $A$ is the interval $[-2,2.66]$ and if $B$ is the interval $(0,10)$, then $A \cup B=(0,2.66]$. True False\\
(d) If $A$ is the interval $[-2,2.66]$ and if $B$ is the interval $(0,10)$, then $A \cap B=(0,2.66]$. True False

\section*{Exam 1 - Form A}
Problem 14. Consider the following system of inequalities.

$$\begin{array}{r}
-2 x+y \geq 9 \\
2 x+y<3
\end{array}$$

Below you are given three ordered pairs written in the form $(x, y)$. CIRCLE the ordered pairs that ARE solutions to the system of inequalities given above.

$$(-5,2) \quad(-2,5) \quad(0,0)$$

You may use the space below to show your work.


\end{document}